\documentclass{article}
\usepackage[left=3cm, right=3cm, top=2cm]{geometry}

\usepackage[utf8]{inputenc}
\usepackage{amsmath}
\usepackage{booktabs}
\usepackage{graphicx}
\usepackage{float}
\usepackage[obeyspaces]{url}

\usepackage{hyperref}
\hypersetup{
    colorlinks=true,
    linkcolor=blue,
    filecolor=magenta,
    urlcolor=cyan,
}

% Image if figs/<file> exists, otherwise an empty box of the given width.
\newcommand{\img}[2]{% file, width
  \IfFileExists{figs/#1}
    {\includegraphics[width=#2,height=0.3\textheight,keepaspectratio]{figs/#1}}
    {\fbox{\parbox[c][5cm][c]{#2}{\centering Image: \texttt{figs/\detokenize{#1}}}}}}

% Figure slot with one image.
\newcommand{\pic}[4][0.7\textwidth]{% [width], file, caption, label
  \begin{figure}[H]
    \centering
    \img{#2}{#1}
    \caption{#3}
    \label{#4}
  \end{figure}}

% Figure slot with two images side by side.
\newcommand{\picpair}[4]{% file1, file2, caption, label
  \begin{figure}[H]
    \centering
    \img{#1}{0.45\textwidth}\hspace{0.5cm}\img{#2}{0.45\textwidth}
    \caption{#3}
    \label{#4}
  \end{figure}}

\title{DEIS Lab 2: Basic Robotics 2  \\ DT8007 VT26, Halmstad University}
\author{Binay Kumar Sah,  binsah25@student.hh.se, Master's in Embedded and Intelligence System
        \and
        Kabir Tamari, kamtam25@student.hh.se, TAEIS 25}
\date{1 October 2026}

\begin{document}

\maketitle

\section*{Introduction}
In this lab we programmed an Arduino UNO R4 WiFi, did some basic hardware
prototyping, flew a Ryze Tello drone (manually and from Python), and got a
SparkFun RedBot ground robot to move.

% ---------------- (a) Arduino ------------------------
\section{Arduino}

\subsection{Getting started}
We installed the UNO R4 board package in the Arduino IDE, selected the board and
port, and uploaded the Blink example. We ran it on a classic Arduino Uno as well
as on the UNO R4 WiFi.

\picpair{uno-blk-off.jpg}{uno-blk-on.jpg}{Blink example on an Arduino Uno: LED L
off (left) and on (right).}{fig:blink}

\pic{unor4-2.jpg}{Sketch running on the Arduino UNO R4 WiFi (LED L on).}{fig:blink-r4}

\subsection{Connecting to WiFi}
We uploaded the \texttt{SimpleWebServerWiFi} example with the course router's
SSID and password, opened the IP address from the Serial Monitor in a browser,
and turned the LED on and off (\texttt{/H} and \texttt{/L}).

\pic{unor4.jpg}{LED on pin 13 turned on from the browser over WiFi. The terminal
on the left shows the HTTP requests received by the board.}{fig:wifi}

% ---------------- (b) Prototyping ------------------------
\section{Prototyping}
Together with the supervisor we did simple tasks with a multimeter, soldering
iron and drill. We used the multimeter to check batteries, and soldered LEDs and
wires onto a perfboard held in a helping-hands stand.

\pic{btrcheck.jpt.jpg}{Multimeter set up for checking batteries.}{fig:proto-multimeter}

\pic{solder.jpg}{Perfboard with an LED in the helping hands, next to the
soldering station.}{fig:proto}

\pic{solder2.jpg}{Solder joints on the perfboard.}{fig:proto-joints}

% ---------------- (c) Drone ------------------------
\section{Drone}

\subsection{Manual flying}
Using the Ryze Tello app we took off, rotated left and right, flew up and down,
forward and backward, strafed left and right, and landed.

\pic{dj-manual.jpg}{Tello hovering in the lab during manual flight.}{fig:drone-manual}

\subsection{Control via Python}
\pic[\textwidth]{cmd-dj.jpg}{Controlling the drone from Python. The terminal shows the
commands sent, the responses, and the saved camera frames.}{fig:drone-python}

We then added a simple autonomous sequence that only rotates: take off, wait
1\,s, rotate 90$^\circ$ clockwise, wait 1\,s, rotate 90$^\circ$ counter-clockwise, wait 1\,s,
and land.

\pic[\textwidth]{telo-py.jpg}{Output of the autonomous sequence (take off, rotate,
land).}{fig:drone-auto}

% ---------------- (d) Ground robot ------------------------
\section{Ground robot}
We assembled the RedBot (main board, IR line sensors, wheel encoders, alkaline
battery pack) and worked through the SparkFun tutorials 1--3 (basic motions),
6 (line following) and 7 (encoders)\footnote{\url{https://learn.sparkfun.com/tutorials/experiment-guide-for-redbot-with-shadow-chassis}}.

\paragraph{Comment.} We did not have black tape to make a line-following
track, so we connected the RedBot to the Arduino UNO R4 WiFi over serial and
drove it from a browser instead. Commands go from WiFi to the UNO and from the
UNO to the RedBot:
\[
  \text{Browser} \xrightarrow{\text{WiFi}} \text{UNO R4 WiFi}
  \xrightarrow{\text{serial RX/TX}} \text{RedBot}
\]
The UNO TX (pin 1) is wired to the RedBot RX (A1), the UNO RX (pin 0) to the
RedBot TX (A0), and the grounds are connected; the link runs at 9600 baud. The
UNO creates its own WiFi network and serves a page with W/A/S/D controls. Each
key press reaches the UNO as an HTTP request, and the UNO passes the letter on
to the RedBot over serial. The RedBot drives the motors and replies with its IR
line-sensor and encoder readings, which the UNO shows live on the page. This is
how we tested the basic motions, the line sensors and the encoders without a
track. As a safety measure, the RedBot stops if no command arrives for 500\,ms.

\pic{redbot.jpg}{RedBot with the UNO R4 WiFi mounted on top, wired over serial
RX/TX, and powered by a power bank.}{fig:redbot}

\pic[0.5\textwidth]{redbot-contorl.png}{Web control page served by the UNO R4 WiFi, with live IR
(left/centre/right) and encoder readings.}{fig:redbot-web}

% ---------------- Contribution ------------------------
\section*{Contribution}
\begin{table}[H]
  \centering
  \begin{tabular}{lcc}
    \toprule
    Activity & Binay Kumar Sah & Kabir Tamari \\
    \midrule
    (a) Arduino      & 50\% & 50\% \\
    (b) Prototyping  & 50\% & 50\% \\
    (c) Drone        & 80\% & 20\% \\
    (d) Ground robot & 20\% & 80\% \\
    Report           & 50\% & 50\% \\
    \bottomrule
  \end{tabular}
\end{table}

\end{document}
